% ============================================================
%  第6章  总结与展望
% ============================================================
\chapter{总结与展望}

\section{工作总结}

本文以脑卒中后视觉功能障碍的居家康复需求为牵引，设计并研制了一套基于STM32H723VGT6微控制器的低成本头戴式眼部康复训练装置。研究路线从临床需求分析切入，依次完成了系统总体方案论证、硬件平台选型与电路搭建、嵌入式软件架构设计与算法实现、以及系统级功能验证的闭合链条。以下对各层面主要工作进行归纳。

系统方案层面，遵循"感知—决策—执行—反馈"的闭环控制范式，定义了覆盖五种临床视觉功能的康复训练模式集合。针对脑卒中后视觉障碍的高发生率、既有干预措施普遍存在的"高端设备用不起、低端手段没量化"的两端困境，以及神经可塑性理论所支持的规律性长期训练需求，将系统的设计定位确立为"低成本、便携式、可量化、全语音操作"四个核心维度。


硬件设计层面，以DM-MC02核心板为平台，成功集成了BMI088六轴惯性传感器、SG90双轴舵机云台、650nm红色激光模组、CI1302 AI语音交互模块、WS2812可编程状态指示灯、板载蜂鸣器和患者反馈按键。机械外壳采用了成品自行车头盔，并手动集成了开发板和各模块到头盔上，电源和按键模块集成为一体独立于头戴装置外。舵机云台、语音模块外壳、开发板外壳、电源和按键模块外壳均采用soildworks自行绘制并通过3D打印形成PLA材质的机械结构，具有轻量化、造价低，同时兼顾耐用性的特点。

 软件设计层面，在STM32H7平台完成了完整的前后台嵌入式软件开发——实现了基于互补滤波和纯陀螺仪积分的多策略姿态估计算法，提出并验证了从传感器坐标系到佩戴者头部坐标系的线性重映射方案，建立了100帧滑动窗口的头稳指标计算、代偿性转头自动检测和多级姿态安全评估体系。设计并实现了五状态系统状态机（系统空闲等待状态 → 系统姿态校准状态 → 系统训练执行状态 → 系统反馈评价状态 → 系统安全暂停状态），五种训练模式，音频播报分级评价机制（基于正确率/反应时间/头稳指标），以及3秒连续监测的自动暂停/恢复策略和全局语音紧急中断保护。语音交互系统支持10条可识别命令词和45条TTS播报词，设计了唤醒词区与命令词区的独立编码映射方案。

人因设计层面，围绕老年脑卒中患者的特点实施了多项实用性适配：全部操作流程压缩为语音命令加单次按键的组合，免去了屏幕显示和多层菜单导航的依赖；所有播报语音采用较慢语速并设置了充裕的句间静音间隔以防音频串扰；标定模式通过按键逐次标记视野边界的操作方式，降低了个体间视野范围差异带来的适配难度。应当指出，上述设计策略的定量有效性目前尚缺乏真实患者群体的实测数据支持，有待后续临床验证加以检验。


\section{本设计工作的不足}

本装置作为面向医疗器械应用方向的探索性工程成果，尽管在软硬件集成和系统功能闭环方面已具备一定的完整性，但以临床康复器械的标准来审视，仍然存在若干不容忽视的短板与边界约束。

临床验证环节的缺失是现阶段最核心的差距。装置目前的测试范围局限于开发者自身及数名健康志愿者的功能性调试，尚未在脑卒中后视觉障碍的真实患者人群中开展过任何形式的临床试用或对照实验。五种训练模式所声称的康复收益——对视野缺损程度的缓解、眼球运动灵活性的恢复以及空间感知能力的改善——均缺乏标准临床量表的同步对比数据作为支撑。例如，训练前后的Fugl-Meyer视觉分量表评分、Humphrey视野计检测的定量改善值以及Barthel日常生活能力指数的变化量等关键结局指标，均未经过系统采集与分析。从循证医学的角度出发，装置仍处于工程原型阶段，距离经过规范临床试验检验的康复器械尚有一段实质性距离。

眼动测量能力的空白构成了技术层面的根本性制约。本装置通过头部运动间接推测眼球行为——例如以代偿性转头作为患者"以头动替代眼动"的代理指标。但该间接推断存在不可忽视的盲区：患者完全可能在保持头部纹丝不动的同时并未进行任何有效的眼球运动（即注意力已分散至训练任务之外），系统对这类状态不具备分辨能力。实现对注视点坐标、扫视频率与扫视幅值等直接眼动参数的精确采集，是视觉康复训练走向高质量客观量化的关键技术支点，而本装置目前尚未集成这一能力。

姿态漂移与坐标系的重标定需求也是工程迭代中需正视的问题。尽管陀螺仪零偏采用了800样本采样的均值校准方案，横滚角（roll）依然全程依赖角速度积分加以维持，在持续运行超过十分钟后积累的漂移量将不再可忽略。患者在训练间歇摘下装置后再次佩戴时，横滚角的绝对基准值会发生偏移，重新执行一次完整的校准流程方可恢复坐标精度。除此之外，偏航角（yaw）完全依靠加速度计重力分量反算获得绝对参考，虽能锁定基准方位，但对缓变运动和姿态恢复的动态跟踪响应偏慢。

激光安全性及使用环境方面的约束同样值得正视。装置搭载的650nm Class II级激光在普通室内照明下光点对比度充足，但切换至强光环境（如阳光经窗直射墙面）后光斑的可辨识度急剧下降直至不可见，使装置对使用场景的光照条件形成了隐性依赖。此外，虽然Class II在分类标准上属于安全性较高的等级，在头戴式佩戴的背景下，若机械固定结构出现松动或者佩戴者姿态出现意外偏转，出射光束存在直接或经反射进入使用者或旁观者眼部的潜在风险，构成不可忽略的安全隐患。

机械结构的材质与佩戴工学方面，装置主体框架采用PLA材料通过3D打印制造，虽然赋予了快速原型迭代和低物料成本的灵活性，但PLA的脆性特征使其在反复拆装或意外跌落后容易出现断裂失效。头部载荷分布上，云台驱动系统、激光模组和聚焦训练装置集中布置于头盔前部区域，长期佩戴可能引发鼻梁和额部区域的压迫不适；加之目前的头盔模具尺寸调节范围有限，难以覆盖所有体型的佩戴对象。

语音交互模块也存在若干适用性边界。CI1302在安静室内环境中的离线识别准确度较为可靠，但在多人病房或临街窗口等存在环境噪声叠加的场景下，识别成功概率将出现明显衰减。命令词的维护流程要求通过亚博智慧云平台重新生成固件并经由PC完成烧录，无法在设备运行期间动态增删或编辑命令词及其配对播报内容，命令集的可扩展性因而受到限制。此外，播报语速和句间停顿虽已配置为偏慢的参数版本，但由于不同个体对语音节奏的偏好差异显著，这一固定配置难以满足所有患者的个性化需求。

对合并构音障碍或失语症的患者的适用性方面，装置正常运转的隐含前提是使用者能够清晰稳定地发出较长音节的多字命令词（例如"平稳追踪训练"和"空间忽略训练"）。而对于脑卒中后语言功能同时受损的患者——恰好构成了视觉康复潜在需求中的一个重要亚群——语音命令的识别准确性将显著下降，形成一种"最需要该功能的群体反而最难使用该功能"的矛盾性困局。

\section{未来展望}

推动本装置从工程原型走向可实际部署的临床康复工具，需要在以下几个方向上进行持续深入的迭代与拓展。

临床验证与循证评价应列为当前最紧迫的推进事项。后续工作可争取与具备康复医学科或神经内科资质的合作医院联合启动小样本预试验，将装置投入真实脑卒中后视觉障碍患者的日常训练中。研究设计建议采取自身前后对照：筛选一组同向性偏盲（HH）或单侧空间忽略（HSN）的确诊患者，开展为期2至4周、每日1至2次的结构化训练干预，在干预前后分别使用标准化临床量表评估视野缺损程度、眼球运动能力、日常生活独立性与自我效能感的变化幅度。所获数据既能为装置的针对性优化提供事实依据，也可作为后续大样本、多中心的随机对照试验（RCT）的前期证据储备。

融合眼动追踪功能是弥补当前最根本技术短板的必经路径。未来版本可规划在装置前端或近眼位置集成一块微型视频眼动记录（Video Oculography, VOG）相机模组，借助红外光源照射角膜产生的反射光斑（角膜反射法）和瞳孔中心定位算法，在时间轴上实时解算出视线的水平与垂直坐标。该项升级将赋予装置对固视稳定度、扫视潜伏期以及平滑追踪增益等指标的直接量化能力，从根本上提升训练过程的信息采集深度与反馈精度，同时为后续实现难度自适应调节的高级个性化训练策略奠定感知基础。

惯性导航层面的精度提升可沿两条路线并行推进。其一是增配一颗三轴磁力计，将九轴融合方案引入现有姿态解算管线——磁力计提供的绝对航向基准可以有效压制陀螺仪绕竖直轴的积分漂移趋势。其二是在融合算法上迭代至更系统的估计框架，如扩展卡尔曼滤波（EKF）或误差状态卡尔曼滤波（ESKF），在一个统一的递推贝叶斯架构中完成所有惯性传感器和磁力计的异步在线校正。此外，在头盔与头皮的接触界面处嵌入压力或电容式接触传感器，可在患者摘脱装置的瞬间自动触发训练暂停并发出重校准提醒，减少因佩戴状态切换引入的人为姿态误差。

数据驱动的个性化康复路径是中长期迭代中值得探索的方向。装置在每次训练过程中已持续积累了包括正确率、反应时间序列、扫视模式分布、代偿事件计数和头稳趋势等在内的多维行为时间序列，这些数据实质上构成了每位患者康复轨迹的个体化行为画像。后续可将上述数据经脱敏处理后上传至云端平台，借助轻量级机器学习模型对单患者的恢复进度曲线进行建模与预测，并据此自动调整训练参数的配置（如试次密度、目标空间分布策略和超时判定窗口大小），从而实现从固定预设方案向数据驱动的自适应康复调度机制过渡。

远程康复与云端管理架构的引入将显著拓展装置的部署半径。基于现有的UART调试数据输出链路，可扩展为无线回传模块（如WiFi模块或蓝牙BLE信道），使每位患者的每日训练摘要与异常事件日志能够实时汇入医疗云端的中央数据库。康复治疗师可通过医师端管理面板调阅患者的训练日历、效率趋势图与风险事件列表，再据此远程下发调整后的训练方案参数；患者端则可通过配套的手机应用直观查看个人的成绩变化曲线、历史进步里程碑以及系统自动生成的下一次训练建议。对于交通出行不便、医疗资源分布稀疏的基层和农村地区，远程康复方案在降低长期康复的经济和时间成本方面具有突出的现实价值。

无障碍与人因工程的后续深化应聚焦于患者群体内部的异质性需求，构建分层级的交互通道。针对因合并构音障碍或失语症而无法稳定使用语音命令的患者，可将语音输入路径替换为外接的单键扫描交互模块，或利用装置已有的头部姿态感知能力设计基于头部动作的操控范式——如点头确认、摇头取消、转头逐项浏览模式菜单——从而将交互门槛降至最低限度。机械结构层面，可尝试引入TPU等柔性耗材并通过多材料3D打印工艺制造刚柔结合的外壳组件，以提升佩戴贴合度和长时间使用的舒适耐受水平；同时可增设可调配重滑块来优化头盔前后区域的重心分配关系，缓解前部载荷集中造成的单侧压强积聚。

从工程原型向合规医疗器械的角色跨越是装置走向临床服务终端的终极命题。这一转化需要在本装置的设计全流程中引入《医疗器械监督管理条例》及其配套现行标准的合规要求，例如GB 9706系列医用电气设备安全标准~\cite{gb9706_2020}中关于泄漏电流上限、绝缘防护等级和电磁兼容性的硬件级约束，以及YY/T 0287医疗器械质量管理体系标准~\cite{yyt0287_2017}所规定的过程管控框架。软件层面则需要建立与IEC 62304标准~\cite{iec62304}相符合的全生命周期软件工程与验证流程。伦理审查环节需获得所在机构伦理委员会的正式批准并履行患者的书面知情同意程序。尽管这一规范化路径耗时长、投入高，但对于一项承载社会意义和临床潜在价值的康复装置而言，它构成了从实验室走向患者身边所不可绕过的必经之路。

\section{本章小结}

回顾本课题的全过程，完成了一套面向脑卒中后视觉康复、功能闭环且物料成本可控的头戴式训练装置的完整工程实现。在系统集成、软硬件协同调度以及多模态人机交互接口设计三个维度上，项目达成了阶段性的实践产出。同时，也必须坦率地看到现有工作的边界——临床验证样本的空白、直接眼动测量通道的缺失以及若干工程细节尚需打磨，构成了当前版本与成熟康复器械之间不可回避的差距。康复工程这一交叉领域的价值落点，不应停留在展示单项技术的可行性，而在于把技术切实交到需要它的人手中并产生可被客观度量的改善。本文所呈现的工作，希望能成为这条路径上的一个阶段性锚点，在后续研究中逐步填齐临床证据链和人群适用性数据，推动装置从原理样机形态朝循证支撑、具备规模化部署潜力的居家康复工具迈进。
